% Part: computability
% Chapter: recursive-functions
% Section: non-pr-functions

\documentclass[../../../include/open-logic-section]{subfiles}

\begin{document}

\olfileid{cmp}{rec}{npr}
\olsection{આદિમ પુનરાવર્તી ન હોય તેવાં વિધેયો}

આદિમ પુનરાવર્તી વિધેયોમાં સહજ રીતે સંગણનીય લાગતાં બધાં વિધેયો
સમાઈ જતાં નથી. સહજ રીતે સ્પષ્ટ હોવું જોઈએ કે એક આર્ગ્યુમેન્ટવાળાં
બધાં આદિમ પુનરાવર્તી વિધેયોની યાદી
$f_0$, $f_1$, $f_2$,~\dots{} બનાવી શકાય છે, એવી રીતે કે ઇનપુટ
$y$ પર $f_x$નું મૂલ્ય અસરકારક રીતે ગણી શકાય; બીજા શબ્દોમાં, નીચે
મુજબ વ્યાખ્યાયિત વિધેય $g(x,y)$ સંગણનીય છે:
\[
g(x,y) = f_x(y)
\]
પરંતુ ત્યાર પછી નીચેનું વિધેય પણ સંગણનીય છે:
\begin{eqnarray*}
h(x) & = & g(x,x) + 1 \\
& = & f_x(x) +1.
\end{eqnarray*}
દરેક આદિમ પુનરાવર્તી વિધેય $f_i$ માટે $i$ પર $h$ અને $f_i$નાં
મૂલ્યો જુદાં પડે છે. તેથી $h$ સંગણનીય છે, પરંતુ આદિમ પુનરાવર્તી
નથી; $g$ વિશે પણ એ જ કહી શકાય. આ કૅન્ટૉરની વિકર્ણ-રચનાની
``અસરકારક'' આવૃત્તિ છે.

સંગણનીય છતાં આદિમ પુનરાવર્તી ન હોય તેવાં વિધેયોનાં વધુ સ્પષ્ટ
ઉદાહરણો આપી શકાય. દાખલા તરીકે, $g^n(x)$ સંકેત $n$ વખત $g$
વપરાય એવા $g(g(\dots g(x)))$ને દર્શાવે; અને વિધેયોનો અનુક્રમ
$g_0,g_1,\dots$ આમ વ્યાખ્યાયિત કરીએ:
\begin{eqnarray*}
g_0(x) & = & x+1 \\
g_{n + 1}(x) & = & g_n^x(x)
\end{eqnarray*}
દરેક વિધેય $g_n$ આદિમ પુનરાવર્તી છે, તે ચકાસી શકાય. દરેક
આગળનું વિધેય તેના પહેલાંના કરતાં ઘણું ઝડપથી વધે છે: $g_1(x)$ એ
$2x$ છે, $g_2(x)$ એ $2^x \cdot x$ છે અને $g_3(x)$ આશરે $x$
સંખ્યામાં $2$ની ઘાત-મિનારની જેમ વધે છે. અકરમાન--પેટર વિધેય
મૂળભૂત રીતે $G(x) = g_x(x)$ વિધેય છે, અને તે કોઈ પણ આદિમ
પુનરાવર્તી વિધેય કરતાં વધુ ઝડપથી વધે છે એવું બતાવી શકાય છે.

હવે આદિમ પુનરાવર્તી વિધેયોની પરિગણનાના પ્રશ્ન તરફ પાછા ફરીએ.
યાદ કરો કે દરેક આદિમ પુનરાવર્તી વિધેયને આપણે પ્રતીકાત્મક
સંકેતન આપ્યું છે; તેથી સંકેતનોની પરિગણના કરવી પૂરતી છે. દરેક
સંકેતન $F$ને પ્રાકૃતિક સંખ્યા $\#(F)$ પુનરાવર્તી રીતે આ મુજબ
આપી શકીએ:
\begin{eqnarray*}
\#(0) & = & \langle 0 \rangle \\
\#(S) & = & \langle 1 \rangle \\
\#(\Proj{n}{i}) & = & \langle 2, n, i \rangle \\
\#(\fn{Comp}_{k,l}[H,G_0,\dots,G_{k-1}]) & = & \langle
3,k,l,\#(H),\#(G_0),\dots,\#(G_{k-1}) \rangle \\
\#(\fn{Rec}_l[G,H]) & = & \langle 4, l, \#(G), \#(H) \rangle
\end{eqnarray*}
અહીં પાછલા વિભાગના કોડ વાપરીને સંખ્યાઓના દરેક અનુક્રમને પ્રાકૃતિક
સંખ્યા તરીકે જોવાની હકીકત કામમાં લઈએ છીએ. પરિણામે દરેક કોડને
પ્રાકૃતિક સંખ્યા મળે છે. અલબત્ત, કેટલાક અનુક્રમો (અને તેથી કેટલીક
સંખ્યાઓ) કોઈ સંકેતનને અનુરૂપ નથી. પરંતુ $i$ કોઈ સંકેતનનો કોડ
હોય, તો $f_i$ એ $i$ વડે કોડ થયેલું સંકેતન ધરાવતું એક-આર્ગ્યુમેન્ટવાળું આદિમ
પુનરાવર્તી વિધેય માનીએ; અન્યથા તે હંમેશાં $0$ આપતું વિધેય માનીએ.
આ રીતે એક-આર્ગ્યુમેન્ટવાળાં આદિમ પુનરાવર્તી વિધેયોની સ્પષ્ટ
પરિગણના મળે છે.

(હકીકતમાં, અચળ શૂન્ય વિધેય જેવાં કેટલાંક વિધેયો યાદીમાં એકથી વધુ
વખત આવશે. આ માત્ર આપણા સંકેતીકરણની આડઅસર નથી; અચળ શૂન્ય
વિધેયને એક કરતાં વધુ સંકેતનો મળે છે તેનું પણ આ પરિણામ છે.
પછી આપણે જોઈશું કે આવી પુનરાવૃત્તિઓ સંગણનીય રીતે દૂર કરી શકાતી
નથી. દાખલા તરીકે, આપેલું સંકેતન અચળ શૂન્ય વિધેય દર્શાવે છે કે
નહીં તે નક્કી કરતું કોઈ સંગણનીય વિધેય નથી.)

હવે આપણે $g(x,y)$ને $f_x(y)$ વડે આપેલું માનીએ, જ્યાં $f_x$ હમણાં
વર્ણવેલી પરિગણનાનું વિધેય છે. $g(x,y)$ સંગણનીય છે તે કેવી રીતે
જાણીએ? સહજ રીતે આ સ્પષ્ટ છે: $g(x,y)$ ગણવા પહેલાં $x$નો કોડ
``ઉકેલો'' અને જુઓ કે તે એક-આર્ગ્યુમેન્ટવાળા વિધેયનું સંકેતન છે
કે નહીં. હોય તો તે વિધેયનું ઇનપુટ~$y$ પર મૂલ્ય ગણો.

\begin{tagblock}{TMs}
\begin{digress}
તમને કદાચ ખાતરી થઈ ગઈ હશે કે (થોડી મહેનત કરીને!) આ કામ કરતો
પ્રોગ્રામ (કહો કે Java કે C++માં) લખી શકાય. પછી ચર્ચ--ટ્યુરિંગ
માન્યતાનો આધાર લઈ શકીએ: સહજ અર્થમાં જે કંઈ સંગણનીય હોય, તેને
ટ્યુરિંગ મશીનથી ગણી શકાય.

અલબત્ત $g(x,y)$ સંગણનીય છે તે બતાવવાની વધુ સીધી રીત, તેને ગણતું
ટ્યુરિંગ મશીન સ્પષ્ટ રીતે વર્ણવવાની છે. તેમાં ખાસ કરીને
ચર્ચ--ટ્યુરિંગ માન્યતા અને સહજ સમજનો આધાર લેવાની જરૂર નહીં રહે.
ટૂંક સમયમાં આપણે એટલાં સાધનો વિકસાવીશું કે ટ્યુરિંગ મશીન પર
\emph{અનુકરિત} થઈ શકે તેવા સંગણનના એક નમૂનાનો આધાર લઈને
$g(x,y)$ સંગણનીય છે એવું બતાવી શકીશું: એ નમૂનો પુનરાવર્તી
વિધેયોનો છે.
\end{digress}
\end{tagblock}

\end{document}
